\documentclass[12pt]{article}
\usepackage[margin=1in]{geometry}
\usepackage{amsmath,amssymb,amsthm,algorithm,algorithmic}
\usepackage{enumitem}

\newtheorem{theorem}{Theorem}
\newtheorem{lemma}[theorem]{Lemma}

\title{Problem 94: Almost Equilateral Triangles}
\author{Project Euler Solution}
\date{}

\begin{document}
\maketitle

\section{Problem Statement}
An \emph{almost equilateral triangle} has sides $(a, a, a \pm 1)$ with $a \ge 2$. Find the sum of the perimeters of all such triangles with integral area and perimeter $\le 10^9$.

\section{Mathematical Foundation}

\begin{theorem}[Heron's formula, isosceles case]
An isosceles triangle with sides $(a,a,b)$ has area $A = \frac{b}{4}\sqrt{4a^2 - b^2}$.
\end{theorem}
\begin{proof}
With semi-perimeter $s = (2a+b)/2$, Heron's formula gives
\[
A = \sqrt{s(s-a)^2(s-b)} = \sqrt{\frac{2a+b}{2} \cdot \frac{b^2}{4} \cdot \frac{2a-b}{2}} = \frac{b}{4}\sqrt{4a^2 - b^2}. \qedhere
\]
\end{proof}

\begin{lemma}[Case $b = a+1$]
For $(a,a,a+1)$, integral area requires $(3a+1)(a-1) = m^2$ for some $m \in \mathbb{Z}_{>0}$.
\end{lemma}
\begin{proof}
$A = \frac{a+1}{4}\sqrt{(3a+1)(a-1)}$. For $A \in \mathbb{Z}_{>0}$, the radicand must be a perfect square.
\end{proof}

\begin{lemma}[Case $b = a-1$]
For $(a,a,a-1)$, integral area requires $(3a-1)(a+1) = m^2$ for some $m \in \mathbb{Z}_{>0}$.
\end{lemma}
\begin{proof}
$A = \frac{a-1}{4}\sqrt{(3a-1)(a+1)}$. Same argument as above.
\end{proof}

\begin{theorem}[Reduction to Pell's equation]
Both cases reduce to $X^2 - 3Y^2 = 1$.
\end{theorem}
\begin{proof}
\textbf{Case 1:} $(3a+1)(a-1) = m^2$ gives $3a^2 - 2a - 1 = m^2$, so $(3a-1)^2 - 3m^2 = 4$. Setting $X = (3a-1)/2$, $Y = m/2$ yields $X^2 - 3Y^2 = 1$.

\textbf{Case 2:} $(3a-1)(a+1) = m^2$ gives $3a^2 + 2a - 1 = m^2$, so $(3a+1)^2 - 3m^2 = 4$. Setting $X = (3a+1)/2$, $Y = m/2$ yields $X^2 - 3Y^2 = 1$.
\end{proof}

\begin{theorem}[Pell recurrence]
The fundamental solution is $(X_1, Y_1) = (2, 1)$. All positive solutions are generated by
\[
X_{k+1} = 2X_k + 3Y_k, \qquad Y_{k+1} = X_k + 2Y_k.
\]
\end{theorem}
\begin{proof}
This corresponds to multiplication in $\mathbb{Z}[\sqrt{3}]$: $(X_k + Y_k\sqrt{3})(2 + \sqrt{3}) = X_{k+1} + Y_{k+1}\sqrt{3}$. Since the norm $N(X + Y\sqrt{3}) = X^2 - 3Y^2$ is multiplicative:
\[
X_{k+1}^2 - 3Y_{k+1}^2 = (X_k^2 - 3Y_k^2)(2^2 - 3 \cdot 1^2) = 1 \cdot 1 = 1. \qedhere
\]
\end{proof}

\begin{theorem}[Parameter recovery]
From Pell solution $(X_k, Y_k)$:
\begin{itemize}
    \item $X_k \equiv 1 \pmod{3}$: $a = (2X_k+1)/3$, perimeter $= 3a+1$ (Case~1).
    \item $X_k \equiv 2 \pmod{3}$: $a = (2X_k-1)/3$, perimeter $= 3a-1$ (Case~2).
\end{itemize}
The residues $X_k \bmod 3$ alternate as $2, 1, 2, 1, \ldots$
\end{theorem}
\begin{proof}
Inversion of the substitution gives the formulas for~$a$. For alternation: $X_1 = 2 \equiv 2 \pmod{3}$. Since $X_{k+1} \equiv 2X_k \pmod{3}$ and $2^2 \equiv 1 \pmod{3}$, the residues alternate.
\end{proof}

\section{Algorithm}

\begin{algorithm}[H]
\caption{Sum of perimeters of almost equilateral triangles}
\begin{algorithmic}[1]
\REQUIRE Perimeter limit $L$
\STATE $\text{total} \gets 0$, $X \gets 2$, $Y \gets 1$
\LOOP
    \IF{$X \bmod 3 = 1$}
        \STATE $a \gets (2X+1)/3$, $P \gets 3a+1$
    \ELSE
        \STATE $a \gets (2X-1)/3$, $P \gets 3a-1$
    \ENDIF
    \IF{$P > L$} \STATE \textbf{break} \ENDIF
    \IF{$a \ge 2$} \STATE $\text{total} \gets \text{total} + P$ \ENDIF
    \STATE $(X, Y) \gets (2X + 3Y,\; X + 2Y)$
\ENDLOOP
\RETURN total
\end{algorithmic}
\end{algorithm}

\section{Complexity Analysis}

\textbf{Time:} $X_k \sim (2+\sqrt{3})^k$, so the number of iterations is $O(\log L)$. Each iteration is $O(1)$.

\textbf{Space:} $O(1)$.

\section{Answer}

\[
\boxed{518408346}
\]

\end{document}
